\section{Optimal recovery from finite Mellin information}
\label{sec:minimax-recovery}

This section is deterministic.  It states exactly what follows from a
finite list of Mellin samples or jets together with a total-variation
bound, before any special property of the von Mangoldt coefficients is
used.

\subsection{The exact minimax radius}

Let $K\subset\R$ be compact and let $\cM(K)$ denote the complex finite
Borel measures on $K$, with total-variation norm.  For
$\phi_1,\ldots,\phi_J,V\in C(K)$ define
\begin{equation}\label{eq:information-map}
 \cN\mu=\left(\int_K\phi_j\,d\mu\right)_{j=1}^J,
 \qquad
 L_V\mu=\int_K V\,d\mu.
\end{equation}
An algorithm is any map $\Phi:\C^J\to\C$; no linearity or continuity is
assumed.

\begin{theorem}[Exact finite-information minimax identity]
\label{thm:exact-minimax}
For every $R>0$,
\begin{align}
 &\inf_{\Phi:\C^J\to\C}
 \sup_{\|\mu\|_{\TV}\leq R}
 \left|L_V\mu-\Phi(\cN\mu)\right|
 \notag\\
 &\hspace{5em}=
 R\,\dist_{C(K)}\!\left(V,
 \operatorname{span}\{\phi_1,\ldots,\phi_J\}\right).
 \label{eq:exact-minimax}
\end{align}
In particular, nonlinear reconstruction cannot improve on the best linear
reconstruction.
\end{theorem}

\begin{proof}
For $c=(c_j)\in\C^J$, the linear estimator
$\Phi_c(y)=\sum_jc_jy_j$ has error at most
\[
 R\left\|V-\sum_jc_j\phi_j\right\|_{C(K)}.
\]
Taking the infimum over $c$ proves the upper bound.

For the reverse inequality, fix any algorithm $\Phi$ and restrict to data
$\cN\mu=0$.  The two measures $\mu$ and $-\mu$ produce the same data, and
\[
 \max\{|L_V\mu-\Phi(0)|,|-L_V\mu-\Phi(0)|\}\geq|L_V\mu|.
\]
Put $\cE_0=\operatorname{span}\{\phi_1,\ldots,\phi_J\}$.  Its annihilator
in $\cM(K)=C(K)^*$ is
$\cE_0^\perp=\ker\cN$.  The Hahn--Banach distance identity gives
\begin{align*}
 \sup_{\substack{\mu\in\ker\cN\\\|\mu\|_{\TV}\leq1}}|L_V\mu|
 &=\dist_{C(K)}(V,\cE_0).
\end{align*}
Multiplication by $R$ completes the proof.  This is the scalar,
centrally-symmetric form of the classical optimal-recovery principle
\citep{MicchelliRivlin1977}.
\end{proof}

Take sample points $t_1,\ldots,t_q\in\R$ and multiplicities
$s_1,\ldots,s_q\geq1$.  The observations
\begin{equation}\label{eq:Mellin-jet-data}
 \widehat\mu^{(k)}(t_\ell)
 =\int_K(iu)^ke^{it_\ell u}\,d\mu(u),
 \qquad 0\leq k<s_\ell,
\end{equation}
correspond to the information space
\begin{equation}\label{eq:exponential-polynomial-space}
 \cE(\boldsymbol t,\boldsymbol s)
 =\operatorname{span}\{u^ke^{it_\ell u}:
 1\leq\ell\leq q,\ 0\leq k<s_\ell\}.
\end{equation}
Consequently \cref{thm:exact-minimax} becomes
\begin{equation}\label{eq:Mellin-minimax-distance}
 \inf_\Phi\sup_{\|\mu\|_{\TV}\leq R}
 \left|L_V\mu-\Phi((\widehat\mu^{(k)}(t_\ell))_{\ell,k})\right|
 =R\,\dist_{C(K)}(V,\cE(\boldsymbol t,\boldsymbol s)).
\end{equation}
If only the jets $0,\ldots,K_0$ at $t=0$ are known, the recovery space is
the polynomial space $\cP_{K_0}$.

\begin{proposition}[Exact noise--resolution tradeoff]
\label{prop:noisy-minimax}
Suppose the observed coordinates are
\(
 y_j=\int\phi_j\,d\mu+e_j
\)
with $\|\mu\|_{\TV}\leq R$, $|e_j|\leq\eps_j$, and every $\eps_j>0$.
Then the worst-case
minimax error is
\begin{equation}\label{eq:noisy-minimax}
 \inf_{c\in\C^J}
 \left\{
 R\left\|V-\sum_{j=1}^Jc_j\phi_j\right\|_\infty
 +\sum_{j=1}^J\eps_j|c_j|
 \right\}.
\end{equation}
\end{proposition}

\begin{proof}
Apply the proof of \cref{thm:exact-minimax} to the product variable
$(\mu,e)$ with observation $\cN\mu+e$ and target $L_V\mu$.  The uncertainty
set is the unit ball of
\[
 \max\left\{\|\mu\|_{\TV}/R,
             \max_j|e_j|/\eps_j\right\},
\]
whose dual norm is
$R\|f\|_\infty+\sum_j\eps_j|c_j|$.  The distance from $(V,0)$ to the
range of the adjoint observation map is exactly
\eqref{eq:noisy-minimax}.
\end{proof}

\subsection{A fixed finite data set becomes asymptotically trivial}

The log-scale range of factor pairs with $mn\asymp X$ has length
$2\log X+O_W(1)$.  On an expanding range, every fixed finite set of
frequencies and jet orders eventually performs no better, in the minimax
sense, than the constant midpoint approximation.

\begin{theorem}[The fixed-data midpoint limit]
\label{thm:midpoint-limit}
Let $V\in C_c(\R)$ satisfy $0\leq V\leq M$ and $\max V=M>0$.  Fix a finite
set of real frequencies and jet multiplicities, independent of $L$, and
let $\cE$ be the associated space \eqref{eq:exponential-polynomial-space}.
Assume $1\in\cE$.  Then
\begin{equation}\label{eq:midpoint-limit}
 \lim_{L\to\infty}\dist_{C([-L,L])}(V,\cE)=\frac{M}{2}.
\end{equation}
\end{theorem}

\begin{proof}
The constant $M/2$ gives the upper bound.  Suppose, to the contrary, that
there are $\eta<M/2$, numbers $L_j\to\infty$, and $q_j\in\cE$ with
\(
 \|V-q_j\|_{C([-L_j,L_j])}\leq\eta
\).
On any fixed nondegenerate interval the restriction map on the
finite-dimensional space $\cE$ is injective, so norm equivalence makes the
coefficient vectors of $q_j$ bounded.  A subsequence converges in $\cE$ to
an exponential polynomial $q$, and
\begin{equation}\label{eq:global-midpoint-approximant}
 \|V-q\|_{C(\R)}\leq\eta.
\end{equation}

In particular $q$ is bounded on $\R$.  A bounded exponential polynomial
$\sum_\ell P_\ell(u)e^{it_\ell u}$ with distinct real $t_\ell$ has all
$P_\ell$ constant: divide by the largest power of $u$ and use the positive
mean square of a nonzero trigonometric polynomial.  Hence $q$ is a finite
trigonometric polynomial.  Such a function is uniformly almost periodic;
there are arbitrarily large $\tau$ for which
\begin{equation}\label{eq:almost-periodic-return}
 \sup_{u\in\R}|q(u+\tau)-q(u)|<\delta
\end{equation}
for any $\delta>0$; see \citep{Corduneanu1989}.

Choose $u_0$ with $V(u_0)=M$ and then choose an almost period $\tau$ so
large that $u_0+\tau\notin\supp V$.  From
\eqref{eq:global-midpoint-approximant} and
\eqref{eq:almost-periodic-return},
\[
 M-\eta\leq|q(u_0)|
 \leq|q(u_0+\tau)|+\delta\leq\eta+\delta.
\]
Letting $\delta\downarrow0$ contradicts $\eta<M/2$.
\end{proof}

\begin{remark}[Fixed means fixed]
The frequencies and jet multiplicities in
\cref{thm:midpoint-limit} do not vary with $L$.  If they are allowed to
depend on $X$, the exact distance formula
\eqref{eq:Mellin-minimax-distance} remains true at each scale, but the
limit \eqref{eq:midpoint-limit} does not follow from this argument.
\end{remark}

\subsection{How many zero-frequency jets are needed?}

The next result gives the correct information-level order for an interior
window.  It does not say that the required arithmetic derivatives can be
estimated.

\begin{theorem}[Interior jet complexity]
\label{thm:jet-complexity}
Let $0<w\leq L/2$, let $0\leq V\leq1$ on $[-L,L]$, and suppose
$V(0)=1$ and $V(w)=0$.  If a polynomial $p$ of degree at most $K$ satisfies
\begin{equation}\label{eq:jet-approximation-error}
 \|V-p\|_{C([-L,L])}\leq\eps<\frac12,
\end{equation}
then
\begin{equation}\label{eq:jet-complexity-lower}
 K\geq
 \frac{\sqrt3}{2}\frac{1-2\eps}{1+\eps}\frac{L}{w}.
\end{equation}
Conversely, if $V$ is Lipschitz with Lipschitz constant $O(1/w)$, Jackson's
theorem gives a polynomial of degree $K$ with
\begin{equation}\label{eq:jet-complexity-upper}
 \|V-p_K\|_{C([-L,L])}\ll \min\left\{1,\frac{L}{Kw}\right\}.
\end{equation}
\end{theorem}

\begin{proof}
The mean-value theorem and \eqref{eq:jet-approximation-error} give a
$\xi\in(0,w)$ such that
\[
 |p'(\xi)|\geq\frac{1-2\eps}{w}.
\]
Bernstein's interior inequality on $[-L,L]$ gives
\[
 |p'(\xi)|\leq
 \frac{K}{\sqrt{L^2-\xi^2}}\|p\|_\infty
 \leq\frac{2K(1+\eps)}{\sqrt3L}.
\]
This proves \eqref{eq:jet-complexity-lower}.  After rescaling
$[-L,L]$ to $[-1,1]$, Jackson's inequality bounds the best degree-$K$
error by a constant times the modulus of continuity at scale $1/K$,
which is $O(L/(Kw))$ under the stated Lipschitz hypothesis; see
\citep{Lorentz1986}.
\end{proof}

For factor scales, one may take $L=\log X+O_W(1)$.  Thus a fixed-width
interior window requires $K\asymp\log X$ zero-frequency jets for constant
accuracy at the deterministic recovery level.  Whether arithmetic can
control that many jets is a separate and substantially harder question.

\subsection{What a continuum band would provide}

In this subsection assume additionally that
$V\in C_c^k(\R)$ for some $k\geq2$.
Let
\(
 \mathscr F V(t)=\int_\R V(u)e^{-itu}\,du
\)
and
\(
 \widehat\mu(t)=\int_\R e^{itu}\,d\mu(u)
\).
Fourier inversion gives
\begin{equation}\label{eq:continuum-band-inversion}
 \int_\R V\,d\mu
 =\frac1{2\pi}\int_\R \mathscr F V(t)\widehat\mu(t)\,dt.
\end{equation}
If only $|t|\leq T$ is controlled, the deterministic truncation error is
\begin{equation}\label{eq:continuum-band-error}
 \leq\frac{\|\mu\|_{\TV}}{2\pi}
 \int_{|t|>T}|\mathscr F V(t)|\,dt
 \ll_k\|\mu\|_{\TV}\|V^{(k)}\|_1T^{1-k}
\end{equation}
for $k\geq2$.  Also, for
$A_V(s)=\int V(u-s)\,d\mu(u)$, Plancherel yields
\begin{equation}\label{eq:scale-center-Plancherel}
 \int_\R|A_V(s)|^2\,ds
 =\frac1{2\pi}\int_\R
 |\mathscr F V(t)|^2|\widehat\mu(t)|^2\,dt.
\end{equation}
Thus a genuine continuum-band estimate can imply averaged scale
localization.  Isolated finite samples cannot substitute for such an
estimate without additional regularity information.
